\documentclass[10pt,a4paper]{amsart}
\usepackage{amsmath,amssymb,amsthm,graphicx}
% Retained preamble for the Stacks Project chapter {CH} (shared project preamble).
% The chapter file pinned-original-source.tex only carries \input{{preamble}} on its first line,
% so no preamble text lives inside the chapter itself.  The Stacks Project uses one single shared
% preamble for every chapter: the snapshot stages it three times (sources/stacks-chow/preamble.tex,
% sources/stacks-etale-cohomology/preamble.tex and focused-seeds/AGP/sources/stacks-snapshot-20261011/
% preamble.tex) and those copies are identical in size.  The block below is copied mechanically,
% line for line, from sources/stacks-chow/preamble.tex lines 14-18 and 154-217: the Xy-pic setup and
% the native theorem declarations together with the mathematical macros.  Nothing was added, removed
% or redefined; the wrapper class and the amsmath/amssymb/amsthm/graphicx base packages are supplied
% by the assembler.
% For commutative diagrams we use Xy-pic
% Chapter: stacks-spaces-more-morphisms.tex -- More on Morphisms of Spaces.
\usepackage[all]{xy}

% We use 2cell for 2-commutative diagrams.
\xyoption{2cell}

% Theorem environments.
%
\theoremstyle{plain}
\newtheorem{theorem}[subsection]{Theorem}
\newtheorem{proposition}[subsection]{Proposition}
\newtheorem{lemma}[subsection]{Lemma}

\theoremstyle{definition}
\newtheorem{definition}[subsection]{Definition}
\newtheorem{example}[subsection]{Example}
\newtheorem{exercise}[subsection]{Exercise}
\newtheorem{situation}[subsection]{Situation}

\theoremstyle{remark}
\newtheorem{remark}[subsection]{Remark}
\newtheorem{remarks}[subsection]{Remarks}

\numberwithin{equation}{subsection}

% Macros
%
\def\lim{\mathop{\mathrm{lim}}\nolimits}
\def\colim{\mathop{\mathrm{colim}}\nolimits}
\def\Spec{\mathop{\mathrm{Spec}}}
\def\Hom{\mathop{\mathrm{Hom}}\nolimits}
\def\Ext{\mathop{\mathrm{Ext}}\nolimits}
\def\SheafHom{\mathop{\mathcal{H}\!\mathit{om}}\nolimits}
\def\SheafExt{\mathop{\mathcal{E}\!\mathit{xt}}\nolimits}
\def\Sch{\mathit{Sch}}
\def\Mor{\mathop{\mathrm{Mor}}\nolimits}
\def\Ob{\mathop{\mathrm{Ob}}\nolimits}
\def\Sh{\mathop{\mathit{Sh}}\nolimits}
\def\NL{\mathop{N\!L}\nolimits}
\def\CH{\mathop{\mathrm{CH}}\nolimits}
\def\proetale{{pro\text{-}\acute{e}tale}}
\def\etale{{\acute{e}tale}}
\def\QCoh{\mathit{QCoh}}
\def\Ker{\mathop{\mathrm{Ker}}}
\def\Im{\mathop{\mathrm{Im}}}
\def\Coker{\mathop{\mathrm{Coker}}}
\def\Coim{\mathop{\mathrm{Coim}}}

% Boxtimes
%
\DeclareMathSymbol{\boxtimes}{\mathbin}{AMSa}{"02}

%
% Macros for moduli stacks/spaces
%
\def\QCohstack{\mathcal{QC}\!\mathit{oh}}
\def\Cohstack{\mathcal{C}\!\mathit{oh}}
\def\Spacesstack{\mathcal{S}\!\mathit{paces}}
\def\Quotfunctor{\mathrm{Quot}}
\def\Hilbfunctor{\mathrm{Hilb}}
\def\Curvesstack{\mathcal{C}\!\mathit{urves}}
\def\Polarizedstack{\mathcal{P}\!\mathit{olarized}}
\def\Complexesstack{\mathcal{C}\!\mathit{omplexes}}
% \Pic is the operator that assigns to X its picard group, usage \Pic(X)
% \Picardstack_{X/B} denotes the Picard stack of X over B
% \Picardfunctor_{X/B} denotes the Picard functor of X over B
\def\Pic{\mathop{\mathrm{Pic}}\nolimits}
\def\Picardstack{\mathcal{P}\!\mathit{ic}}
\def\Picardfunctor{\mathrm{Pic}}
\def\Deformationcategory{\mathcal{D}\!\mathit{ef}}
\newenvironment{reference}{}{}
\newenvironment{slogan}{}{}
\newenvironment{history}{}{}
\usepackage{cleveref}
\crefname{originallemma}{original lemma}{original lemmas}
\Crefname{originallemma}{Original lemma}{Original lemmas}
\makeatletter
\newlabel{lemma-flat-on-fibres-at-point}{{lemma-flat-on-fibres-at-point}{source}{}{sourcepointer-1}{}}
\newlabel{lemma-flat-on-fibres-at-point@cref}{{[originallemma][0][]lemma-flat-on-fibres-at-point}{[1][0][]source}{}{}{}}
\makeatother
\crefname{originallemma}{original lemma}{original lemmas}
\Crefname{originallemma}{Original lemma}{Original lemmas}
\makeatletter
\newlabel{more-morphisms-lemma-flat-and-free-at-point-fibre}{{more-morphisms-lemma-flat-and-free-at-point-fibre}{source}{}{sourcepointer-2}{}}
\newlabel{more-morphisms-lemma-flat-and-free-at-point-fibre@cref}{{[originallemma][0][]more-morphisms-lemma-flat-and-free-at-point-fibre}{[1][0][]source}{}{}{}}
\makeatother
\crefname{originallemma}{original lemma}{original lemmas}
\Crefname{originallemma}{Original lemma}{Original lemmas}
\makeatletter
\newlabel{spaces-properties-lemma-describe-etale-local-ring}{{spaces-properties-lemma-describe-etale-local-ring}{source}{}{sourcepointer-3}{}}
\newlabel{spaces-properties-lemma-describe-etale-local-ring@cref}{{[originallemma][0][]spaces-properties-lemma-describe-etale-local-ring}{[1][0][]source}{}{}{}}
\makeatother
\crefname{originallemma}{original lemma}{original lemmas}
\Crefname{originallemma}{Original lemma}{Original lemmas}
\makeatletter
\newlabel{spaces-properties-lemma-stalk-quasi-coherent}{{spaces-properties-lemma-stalk-quasi-coherent}{source}{}{sourcepointer-4}{}}
\newlabel{spaces-properties-lemma-stalk-quasi-coherent@cref}{{[originallemma][0][]spaces-properties-lemma-stalk-quasi-coherent}{[1][0][]source}{}{}{}}
\makeatother
\title{A finitely presented module flat at a point is free there exactly when its fibre module is free}
\author{The Stacks Project Authors}
\date{stacks-project:spaces-more-morphisms.tex:local-snapshot-20261011; GNU FDL 1.2 or later}
\begin{document}
\maketitle
\noindent\textit{Editorial note.} Statement span is lines 5939-5957 of sources/stacks-spaces-more-morphisms/pinned-original-source.tex; proof span is lines 5959-5986, the complete author proof (26 body lines).\par\medskip
\noindent\textit{Editorial note.} Navigation stubs: 4 cited-but-not-included labels are supplied as non-load-bearing pointers with their pinned tags where the snapshot tag index has them.\par\medskip
\noindent\textit{Editorial note.} extra\_packages is empty on purpose: the retained preamble already loads xy with options and a second declaration clashes in the assembler wrapper.\par\medskip
\noindent\textit{Editorial note.} Text is retained verbatim, including the author's own typos and placeholders.\par\medskip
\noindent\textit{Allowed standard prerequisites (editorial).}\begin{itemize}
\item Explicit prerequisite boundary. Finitely presented modules flat over a morphism, the local freeness criterion, and etale neighbourhoods; the descent is retained.
\end{itemize}\medskip
\ifdefined\hypertarget\hypertarget{sourcepointer-1}{}\fi
\noindent\textit{Source reference.} Original lemma lemma-flat-on-fibres-at-point: sources/stacks-spaces-more-morphisms/pinned-original-source.tex, source lines 5632--5662. Cited by the retained statement or proof but not part of the collected spans; supplied only as a navigation pointer. The argument does not depend on its content.\par\medskip
\ifdefined\hypertarget\hypertarget{sourcepointer-2}{}\fi
\noindent\textit{Source reference.} Original lemma more-morphisms-lemma-flat-and-free-at-point-fibre: sources/stacks-spaces-more-morphisms/pinned-original-source.tex, source lines 5986--5986. Cited by the retained statement or proof but not part of the collected spans; supplied only as a navigation pointer. The argument does not depend on its content.\par\medskip
\ifdefined\hypertarget\hypertarget{sourcepointer-3}{}\fi
\noindent\textit{Source reference.} Original lemma spaces-properties-lemma-describe-etale-local-ring: sources/stacks-spaces-more-morphisms/pinned-original-source.tex, source lines 5986--5986. Cited by the retained statement or proof but not part of the collected spans; supplied only as a navigation pointer. The argument does not depend on its content.\par\medskip
\ifdefined\hypertarget\hypertarget{sourcepointer-4}{}\fi
\noindent\textit{Source reference.} Original lemma spaces-properties-lemma-stalk-quasi-coherent: sources/stacks-spaces-more-morphisms/pinned-original-source.tex, source lines 5986--5986. Cited by the retained statement or proof but not part of the collected spans; supplied only as a navigation pointer. The argument does not depend on its content.\par\medskip
\begin{lemma}
\label{lemma-flat-and-free-at-point-fibre}
Let $S$ be a scheme. Let $f : X \to Y$ be a morphism of algebraic spaces
over $S$ which is locally of finite presentation.
Let $\mathcal{F}$ be a finitely presented $\mathcal{O}_X$-module.
Let $x \in |X|$ with image $y \in |Y|$. If $\mathcal{F}$ is flat at $x$
over $Y$, then the following are equivalent
\begin{enumerate}
\item $(\mathcal{F}_{\overline{y}})_{\overline{x}}$ is a flat
$\mathcal{O}_{X_{\overline{y}}, \overline{x}}$-module,
\item $(\mathcal{F}_{\overline{y}})_{\overline{x}}$ is a free
$\mathcal{O}_{X_{\overline{y}}, \overline{x}}$-module,
\item $\mathcal{F}_{\overline{y}}$ is finite free in an
\'etale neighbourhood of $\overline{x}$ in $X_{\overline{y}}$, and
\item $\mathcal{F}$ is finite free in an \'etale neighbourhood of $x$ in $X$.
\end{enumerate}
Here $\overline{x}$ is a geometric point of $X$ lying over $x$
and $\overline{y} = f \circ \overline{x}$.
\end{lemma}

\begin{proof}
Pick a commutative diagram
$$
\xymatrix{
U \ar[d] \ar[r] & V \ar[d] \\
X \ar[r] & Y
}
$$
where $U$ and $V$ are schemes and the vertical arrows are \'etale
such that there is a point $u \in U$ mapping to $x$. Let $v \in V$
be the image of $u$. Applying Lemma \ref{lemma-flat-on-fibres-at-point}
to $\text{id} : X \to X$ over $Y$ we see that (1) translates into
the condition ``$\mathcal{F}|_{U_v}$ is flat over $U_v$ at $u$''.
In other words, (1) is equivalent to $(\mathcal{F}|_{U_v})_u$
being a flat $\mathcal{O}_{U_v, u}$-module.
By the case of schemes (More on Morphisms, Lemma
\ref{more-morphisms-lemma-flat-and-free-at-point-fibre}),
we find that this implies that
$\mathcal{F}|_U$ is finite free in an open neighbourhood
of $u$. In this way we see that (1) implies (4).
The implications (4) $\Rightarrow$ (3) and
(2) $\Rightarrow$ (1) are immediate.
For the implication (3) $\Rightarrow$ (2) use
the description of local rings and stalks in
Properties of Spaces, Lemmas
\ref{spaces-properties-lemma-describe-etale-local-ring} and
\ref{spaces-properties-lemma-stalk-quasi-coherent}.
\end{proof}

\end{document}
